\begin{table}[htbp]\centering
\caption{Coupled constrained economy: full-prefix work at $N=16$}\label{tab:constrained47}
\begin{tabular}{rrlrrrr}\toprule
$T$ & $p$ & Method & Bound & Queries & Seconds & Time range\\\midrule
2 & 1 & Witness & 2.3148 & 11,534 & 7.045 & [7.003,7.070] \\
2 & 1 & Bilinear FVI & 2.7960 & 11,534 & 0.326 & [0.325,0.352] \\
2 & 4 & Witness & 2.3691 & 11,534 & 6.992 & [6.975,7.188] \\
2 & 4 & Bilinear FVI & 2.8288 & 11,534 & 0.324 & [0.323,0.325] \\
3 & 1 & Witness & 3.3212 & 17,301 & 10.519 & [10.480,10.552] \\
3 & 1 & Bilinear FVI & 3.8915 & 17,301 & 0.490 & [0.487,0.500] \\
3 & 4 & Witness & 3.4107 & 17,301 & 10.482 & [10.425,10.538] \\
3 & 4 & Bilinear FVI & 3.9321 & 17,301 & 0.494 & [0.489,0.520] \\
\bottomrule\end{tabular}
\par\smallskip\begin{minipage}{0.98\linewidth}\footnotesize Query counts and clocks include all three rungs, including failed attempts. Seconds are medians over three isolated process observations; the range reports their minimum and maximum. CPU affinity and library thread counts are fixed, but CPU frequency is not controlled.\end{minipage}\end{table}

Both methods attain the target 4 in all twelve services, at the last rung, and neither attains 2, 1 or $1/2$ within this cap. The witness's final bound is smaller in each economic cell, but bilinear fitted-value iteration is faster in all twelve common successful comparisons. These coarse all-state certificates establish an executed feasible construction, not an economically precise solution at the unachieved tolerances. The sufficient-cap formula states what further resources the mathematical guarantee requires without relabeling these failures.
